\pdfoutput=1
\documentclass[11pt]{article}
\usepackage[T1]{fontenc}
\usepackage{lmodern,microtype}
\usepackage[margin=1in]{geometry}
\usepackage{amsmath,amssymb,amsthm,mathtools}
\usepackage{hyperref}
\hypersetup{colorlinks=true,linkcolor=blue,citecolor=blue,urlcolor=blue,
 pdftitle={Finite scalar tests of matched tensor dynamics}}
\newtheorem{theorem}{Theorem}
\newtheorem{proposition}[theorem]{Proposition}
\newcommand{\dd}{\mathrm{d}}
\newcommand{\ii}{\mathrm{i}}
\newcommand{\bq}{\boldsymbol q}
\newcommand{\bp}{\boldsymbol p}
\newcommand{\TT}{\mathrm{TT}}
\newcommand{\tr}{\operatorname{tr}}
\newcommand{\av}[1]{\left\langle #1\right\rangle}
\newcommand{\D}{\mathcal D}
\newcommand{\I}{\mathcal I}
\newcommand{\R}{\mathcal R}

\title{Finite scalar tests of matched tensor dynamics:\\
Gaussian exchange and an Einstein six-point contrast}
\author{}
\date{}
\begin{document}
\maketitle
\begin{abstract}
A linearly driven Gaussian tensor and a real initial-data representation with the same canonical Gaussian covariance and causal response induce exactly the same closed-time-path influence on quantum scalar probes. Thus a nonzero scalar exchange trispectrum alone does not distinguish these two prescriptions, even at finite time and for arbitrary profiles. At the first pure tensor cubic order they differ by a three-response term of the Einstein action. We compute an explicit nonzero connected six-point coefficient on a finite interval, exhibit its realization by finite real scalar readouts, and prove noncancellation for one scalar species after all fifteen Wick pairings. In a distinct three-species protocol, varying the start time and pair energies separates the bulk tensor contribution from a specified family of time-local finite-derivative mixed omissions. For a finite list of fixed mixed templates the cancellation can be exact on finitely many finite-aperture runs. Propagating non-tensor, apparatus, and initial-boundary sectors are matched, not eliminated. The result is a test of particular response prescriptions in perturbative gravity, not a test against all classical models.
\end{abstract}

\section{Setup and scope}
On the finite conformal-time slab \([s,T]\) take a smooth spatially flat expanding background with \(a>0\). Work in reduced physical variables, with \(\gamma_{ij}=2g h_{ij}\), \(g=M_{\mathrm P}^{-1}\), in the spatial exponential metric \(a^2 e^\gamma\). Set \(\hbar=1\). Spatial indices use the Euclidean metric. For \(\bq\ne0\) let
\[
 P_{ij}=\delta_{ij}-\hat q_i\hat q_j,\qquad
 \Lambda_{ij,kl}=\tfrac12(P_{ik}P_{jl}+P_{il}P_{jk}-P_{ij}P_{kl}).
\]
This is the unit orthogonal projector on symmetric transverse-traceless tensors; there is no extra sum or normalization of polarizations in its use below. In the zero-reference-density tensor limit
\begin{align}
 S_2[h]&=\tfrac12\int \dd\eta\,\dd^3x\,a^2\big[(h_{ij}')^2-(\partial_k h_{ij})^2\big],
 &
 D_q&=\partial_\eta (a^2\partial_\eta)+a^2q^2 .
 \label{eq:S2}
\end{align}
The scalar multiplier of the retarded tensor inverse is \(G_R(q;\eta,v)\), with \(G_R(q;v,v)=0\) and \(a(v)^2\partial_\eta G_R(q;v^+,v)=1\). \(G_A=G_R^t\) in the coordinate pairing. The free probes are real massless conformal scalars; \(\chi_b=a\phi_b\) are in the common conformal Gaussian vacuum with modes \(e^{-\ii p\eta}/\sqrt{2p}\). Its linear TT coupling is
\begin{equation}
 S_{\phi,1}=\int h_{ij}J^{ij},\qquad
 J^{ij}=g\sum_b \partial_i\chi_b\,\partial_j\chi_b .
 \label{eq:scalar}
\end{equation}
At linear TT order the variation of the conformal improvement has \(\delta R=0\) (the background spatial Ricci tensor is isotropic and the TT variation is traceless and divergence-free). This does not discard its second-order mixed terms.

The experiment is defined with common product preparation of the free scalar and physical tensor sectors at \(s\), not an interacting vacuum condition in the past. Every nontested propagating sector, scalar self-coupling, apparatus response, and initial and endpoint calibration is the same between the prescriptions. Section~\ref{sec:null} permits a specified exception for time-local mixed contacts. We use the same declared tensor coordinate and scalar readouts in both prescriptions. The quantum scalar current algebra is never replaced by a classical mean stress. Unless explicitly stated, the expansion is the indicated connected tree coefficient, not a nonperturbatively defined interacting measure. One can use common infrared regulation or states and apertures on which the free covariances exist; all displayed nonzero tests use nonzero transfer momenta.

For a material coordinate calibration, use a weak time label \(X^0\), with \(a^2(X^0)'=\text{constant}\ne0\), and shift-symmetric spatial labels \(X^i=\beta x^i\) with \(\beta\ne0\). The three spatial gradients contribute homogeneous isotropic stress, so they can be included in the background source. Here and below, an expansion support and any other background matter may be part of the common setup. For free minimal reference scalars the linear perturbation of \(\Box_g X^A\) from a TT mode vanishes: \(\delta \sqrt{-g}=0\), \(\delta g^{0\mu}=0\), and \(\partial_i h_{ij}=0\). Hence the scalar readout at common material clock and fixed labels has no linear TT coordinate shift. Higher readout terms and propagating reference modes remain common. This matches coordinate labels, not a state-dependent refitting of geodesic rulers. Nonzero reference stress can be kept weak in the example, cf.\ Section~\ref{sec:packets}.

\section{All-profile Gaussian control}
For two real c-number currents use \(J_r=(J_++J_-)/2\), \(J_a=J_+-J_-\), and make the same change for \(h\). Tensor indices and TT projections are implicit in the pairings. For the centered common Gaussian tensor state let
\[
 N=\tfrac12\av{\{h,h\}},\qquad
 G_R(\eta,\eta')=\ii\theta(\eta-\eta')\av{[h(\eta),h(\eta')]} .
\]
The commutator convention agrees with the retarded inverse above for the source \(+hJ\) in the action. Nonzero common means simply contribute the same linear phase.

\begin{theorem}[Gaussian influence equality]\label{thm:gaussian}
For every pair of admitted source profiles on the finite slab, the quantum tensor trace equals the real Gaussian initial-data average with the same full canonical initial covariance, propagated homogeneously to \(\xi\) and shifted by the retarded drive \(G_R J_r\):
\begin{equation}
 \I_{\rm Q,G}[J_r,J_a]
 =\exp\big(\ii J_aG_RJ_r-\tfrac12 J_aNJ_a\big)
 =\mathbb E\exp\big(\ii J_a[\xi+G_RJ_r]\big).
 \label{eq:gauss}
\end{equation}
In particular the scalar closed-time-path functional built only from these linear tensor exchanges agrees, including the order-\(g^2\) connected exchange four-point function.
\end{theorem}

\begin{proof}
The Wick expansion of the centered tensor trace has only a quadratic cumulant. The \(aa\) tensor covariance vanishes by the forward/backward identities of the two contour branches, the mixed cumulant is the commutator with retarded support, and the \(rr\) covariance is \(N\); this gives the exponent in \eqref{eq:gauss}. Equivalently the bulk quadratic action difference is \(-\int h_a D h_r\) and the source is \(h_aJ_r+h_rJ_a\), with the free state fixing the homogeneous covariance and the final contour trace. A free Gauss\-ian state's symmetrized canonical covariance (including field--momentum cross terms and momentum block) is positive on the real initial test phase space. Its real Gaussian characteristic functional propagated by the same linear dynamics therefore has homogeneous covariance \(N\). The driven solution adds \(G_R J_r\), and evaluating its characteristic functional proves the second equality. These statements can first be made for finitely many Gaussian initial modes and every finite set of admitted linear observables, then extended by consistency of their Gaussian characteristic functions. Substituting the branch currents of the same quantum scalar contour in the common functional proves the scalar claim. Neither a random homogeneous field alone nor a response sourced by an expectation of the scalar current yields this result.
\end{proof}

This is the finite-slab version of the linear influence identity \cite{FeynmanVernon}; it does not construct an autonomous positive hybrid dynamics for arbitrary coupled quantum-classical experiments. The matched retarded term acts on doubled histories of a quantum current. Gaussian initial noises here refer to both canonical blocks, not just a sampled tensor power spectrum.

\section{The pure cubic mismatch}\label{sec:cubic}
The Einstein pure-TT cubic in this tensor coordinate is \cite[Eq.~(37)]{Gao}
\begin{equation}
 S_3[h]=2g\int\dd\eta\,\dd^3x\,a^2
 \left(h_{ik}h_{jl}-\tfrac12h_{ij}h_{kl}\right)\partial_k\partial_l h_{ij}.
 \label{eq:EH3}
\end{equation}
In the exponential coordinate the pure kinetic cubic trace is a trace of a commutator. The linear lapse and shift perturbations sourced by the TT sector vanish; their second-order solutions multiply the vanishing linear constraints and do not enter this pure cubic. This is a statement of pure-sector reduction, not an assumption that the mixed constraints vanish. Any extra pure tensor self-vertex from the background matter is absent or held common except for the optional weak material-ruler term in Section~\ref{sec:packets}.

Compare the full quantum \(S_3\) branch difference with solving \(D h=J_r+S_3'[h]\) perturbatively with the same Gaussian initial draw and averaging \(\exp(\ii J_a h)\) as in \eqref{eq:gauss}. For the comparison in this section keep the complete mixed contour functional the same. The polynomial identity
\[
 S_3[h_r+\tfrac12 h_a]-S_3[h_r-\tfrac12 h_a]
 = S_3'[h_r]\cdot h_a+\tfrac14 S_3[h_a]
\]
implies
\begin{equation}
 \Delta \log \I[J_r,J_a]=\tfrac{\ii}{4} S_3[G_AJ_a]
 \qquad\text{to first order in this pure cubic.}
 \label{eq:pure}
\end{equation}
Here \(\Delta\) is quantum minus the specified trajectory self-dynamics response. To check it directly, the first variation of the averaged trajectory is
\(\ii\mathbb E[e^{\ii J_a h_0}J_aG_R S_3'[h_0]]\),
where \(h_0=\xi+G_RJ_r\). On the Gaussian contour, insertion of \(h_a\) has mean \(G_AJ_a\), and cross contractions between \(h_a\) and \(h_r\) at that same local vertex vanish by \(G_R(v,v)=0\), also with spatial derivatives. The remaining \(h_r\) average is exactly the tilted Gaussian trajectory average. This matches the first term. \(h_a\) has no self-contraction, so the second term gives \eqref{eq:pure} with no \(N\) or tadpole. The argument also applies to the optional weak-reference term below, which has no derivative of \(h\). Common endpoint terms are not varied in this comparison. The absence of multiple-response vertices from a deterministic trajectory truncation is standard \cite{Polkovnikov}. A related triple-retarded mismatch for the inflationary \emph{scalar} cubic was derived by Ballesteros et al.\ \cite[Eq.~(10)]{Ballesteros}; here the three-response tensor vertex is measured through six external scalar fields.

For example the minimal spatial contact \(-g^2\int Q:(h^2)\), \(Q_{ij}=\sum_b\partial_i\chi_b\partial_j\chi_b\), has branch difference
\[
 -g^2\int
 \left[Q_a:(h_r^2+\tfrac14 h_a^2)
       +Q_r:(h_rh_a+h_ah_r)\right].
\]
Matching one-branch coefficients does not match this entire contour contact. The \(Q_a:h_a^2/4\) term is not silently discarded in \eqref{eq:pure}. A comparator that also omits it can differ at the same connected six-point order. Other contour contacts may be present, including conformal-improvement and instantaneous constraint terms.

For the nonlinear expansion one may regulate the real initial trajectories before multiplying fields; there is no assumption that an unsmeared continuum Gaussian sample admits a nonlinear classical solution. Only the displayed finite tree coefficient is compared.

A change only to the initial third tensor cumulant cannot reproduce \eqref{eq:pure} on all c-number currents. Each initial correction leg propagates homogeneously, and tests \(J_{a,b}=D f_b\) with \(f_b\) compactly supported in the slab annihilate all such legs by integration by parts. Yet \(G_A D f_b=f_b\), and the cubic polarization on such triples need not vanish: spatially compact TT tests are obtained by applying the polynomial differential operator with symbol \(q^4\Lambda(\bq)\) to truncations of noncollinear spatial packets chosen in Section~\ref{sec:six}. An expanding common spatial envelope on those packets gives a nonzero leading volume term per volume for the same triple, since derivatives on the envelope are lower order. This argument does not rule out time-dependent non-Gaussian noise, or infer arbitrary-current coverage from a single scalar detector.

\section{Nonzero finite scalar six-point records}\label{sec:six}

\subsection{Contraction and its finite three-species value}
Use \(\chi(\boldsymbol x)=\int \dd^3p\,\chi_{\bp}e^{\ii\bp\cdot\boldsymbol x}/(2\pi)^3\), and initially write the kernels for a late pair \(O_b=\chi_{b,\bp_b}(T)\chi_{b,\bp'_b}(T)\). Below \(p_b=|\bp_b|\) and \(p_b'=|\bp_b'|\). For \(\bq_b=\bp_b+\bp_b'\), the free scalar response is
\begin{align}
 B_b^{ij}(\eta,\boldsymbol x)
 &=\ii\av{[O_b,J_b^{ij}(\eta,\boldsymbol x)]}_0 \notag\\
 &=-g(p_b^ip_b'^j+p_b'^ip_b^j)\,
   \frac{\sin[(p_b+p_b')(T-\eta)]}{2p_bp_b'}e^{-\ii\bq_b\cdot\boldsymbol x}.
 \label{eq:scalarresponse}
\end{align}
Indeed the two connected Wick products of \(O_b J_b\) supply the negative spatial derivative factors and a time phase \(e^{-\ii(p_b+p_b')(T-\eta)}\); in the reverse order the phase is conjugated, and multiplying the difference by \(\ii\) gives \eqref{eq:scalarresponse}.

Let \(V(f_1,f_2,f_3)=[u_1u_2u_3]S_3[\sum_bu_bf_b]\), so \(S_3[h]=V(h,h,h)/6\), and \(L_b=G_A\Lambda B_b\). The connected six-field coefficient from \eqref{eq:pure} is exactly
\begin{equation}
 \Delta \av{O_1O_2O_3}_{\mathrm{conn},3\TT}
 =-\tfrac14 V(L_1,L_2,L_3).
 \label{eq:trivar}
\end{equation}
Each distinct species must attach both its external fields to one current, and the six assignments to the cubic are already in \(V\). The late pair contracts with the difference current as \(-\ii B_b\); the coefficient is \((\ii/4)(-\ii)^3=-1/4\). Disconnected Wick products cannot contribute to this six-field cumulant. The perturbation order is \(g^4\), since the cubic contributes one factor of \(g\) and the three linear current attachments supply the others. This is an isolated tree coefficient; lower-point quantum loop contributions are not compared here.

Choose three equal-length noncollinear equilateral transfers \(\bq_b=q\hat q_b\), with \(\sum_b\bq_b=0\), in a plane normal to \(\hat z\). Let \(n_b=\hat z\times\hat q_b\) and pick each scalar pair
\begin{equation}
 \bp_b=\tfrac12\bq_b+r n_b,\qquad \bp'_b=\tfrac12\bq_b-r n_b,
 \qquad p=\sqrt{q^2/4+r^2},\quad r>0 .
 \label{eq:inplanesplit}
\end{equation}
TT projection of the momentum tensor in \eqref{eq:scalarresponse} gives \(-r^2 M_b\) where \(M_b=n_bn_b^t-\hat z\hat z^t\). Then
\begin{align}
 L_b(v,\boldsymbol x)&=g r^2 M_b e^{-\ii\bq_b\cdot\boldsymbol x} I(q,p;v),
 &
 I(q,p;v)&=\int_v^T\dd\eta\,G_R(q;\eta,v)
                    \frac{\sin[2p(T-\eta)]}{2p^2}.
 \label{eq:leg}
\end{align}
For any permutation \((b,c,d)\), the spatial vertex in the parentheses of \eqref{eq:EH3}, differentiated on the \(d\) leg, gives
\begin{align}
 -(M_b\bq_d)\cdot M_d(M_c\bq_d)&=\tfrac{3}{16}q^2,
 &
 \tfrac12\tr(M_bM_d)\,\bq_d\cdot M_c\bq_d&=\tfrac{15}{32}q^2 .
 \label{eq:angular}
\end{align}
To verify the numbers, \(n_b\cdot n_d=n_c\cdot n_d=-1/2\), \(n_b\cdot\bq_d\) and \(n_c\cdot\bq_d\) have opposite signs and magnitudes \(\sqrt3q/2\), and \(\tr(M_bM_d)=5/4\). Thus the fully symmetrized coefficient of the spatial parentheses is \(63q^2/16\). Removing the common \((2\pi)^3\delta^{(3)}(\sum_b\bq_b)\), we obtain
\begin{equation}
 \boxed{\displaystyle
 \Delta\av{O_1O_2O_3}'_{\mathrm{conn},3\TT}
 =-\tfrac{63}{32}g^4q^2r^6
       \int_s^T\dd v\,a(v)^2\, I(q,p;v)^3 .}
 \label{eq:three}
\end{equation}
This is an exact expression at the specified order. On a sufficiently short nonzero interval,
\(G_R(q;\eta,v)=(\eta-v)/a(v)^2+O((\eta-v)^2)\) uniformly and \(2p(T-\eta)<\pi\). Thus \(I>0\) for \(v<T\). In particular for \(L=T-s\) and fixed \(q,r,a\),
\begin{align}
 I(q,p;v)&=\frac{(T-v)^3}{6p\,a(v)^2}+O((T-v)^4),\\
 \Delta\av{O_1O_2O_3}'_{\mathrm{conn},3\TT}
 &=-\frac{7}{7680}\frac{g^4 q^2r^6}{p^3a(T)^4}L^{10}
   +O(L^{11}).
 \label{eq:shortthree}
\end{align}
No late-time, squeezed, or pole limit is imposed on \eqref{eq:three}.

\subsection{One scalar species without omitting exchange pairings}
For six labeled readings \(X_1,\ldots,X_6\) of one conformal scalar, set
\(\mathcal R_{\ell m}=G_A\Lambda\,\ii\av{[X_\ell X_m,J]}_0\).
The corresponding exact expression is
\begin{equation}
 \Delta C_{6,\mathrm{conn},3\TT}
 =-\tfrac14\sum_{\{(\ell m),(nr),(st)\}}
 V(\mathcal R_{\ell m},\mathcal R_{nr},\mathcal R_{st}).
 \label{eq:allpairs}
\end{equation}
The sum is over all fifteen unordered perfect matchings; the scalar pair contractions and the three tensor-current assignments exhaust the connected tree. In particular, changing from distinct species changes the observable contraction rather than merely its statistics.

For a specific noncancellation, still use the equilateral \(\bq_i\), now with six centers \(\boldsymbol k_{i,\pm}=\bq_i/2\pm b\hat z\), \(0<q\ll b\), of common modulus \(p=\sqrt{b^2+q^2/4}\). For an opposite-sign pair the transfer \(\boldsymbol Q=(\bq_i+\bq_j)/2\) is planar. Its response polarization at size \(b^2\) is \(E(\boldsymbol Q)=\hat z\hat z^t-n(\boldsymbol Q)n(\boldsymbol Q)^t\), with relative correction \(O(q/b)\). In fact the two transverse momenta are \(\pm(b\hat z+u)\) for a planar \(u\perp\boldsymbol Q\) with \(|u|=O(q)\). Their bilinear has an \(O(bq)\) mixed term and a diagonal correction \(O(q^2)\) before TT projection, which proves the bound.

The six opposite-sign pairings are permutations of three minus labels relative to the plus labels. The identity gives the equilateral triad of length \(q\); the two 3-cycles give reversed equilateral triads of length \(q/2\), since \((\bq_i+\bq_j)/2=-\bq_k/2\) for \(i\ne j\). The three transpositions give collinear nonzero transfers, whose spatial contractions vanish exactly by TT transversality. The other three have the sign opposite to \eqref{eq:angular}, because their leading \(E=-M\). Adding the six leading symmetrized coefficients from the parentheses of \eqref{eq:EH3} gives
\begin{equation}
 -\frac{63}{16}\left(q^2+2(q/2)^2\right)=-\frac{189}{32}q^2.
 \label{eq:anglesingle}
\end{equation}
Each of the nine remaining matchings contains one \(++\) and one \(--\) pair. Relative to its own large transfer, each of these two pairs has transverse momentum only \(O(q)\), by subtracting half the pair transfer. Therefore the two TT stresses are \(O(q^2)\) rather than \(O(b^2)\). The two derivatives at the cubic cost at most \(O(b^2)\) rather than \(O(q^2)\), leaving a relative \(O((q/b)^2)\) bound on their summed contribution. All transfers are nonzero.

Write \(a_*=a(T)\), take \(bL\ll1\) and make the relative variation of \(a\) across the slab small. On these finite ranges the standard Volterra equation for \(G_R\) yields its short-time term \((\eta-v)/a(v)^2\) with uniform relative error controlled by \(L\sup|a'/a|\) and \(\max|\boldsymbol Q|^2L^2\). The leading opposite-sign legs are \(g b^2 E(T-v)^3/(6p\,a_*^2)\). Using \eqref{eq:allpairs} and \eqref{eq:anglesingle} gives
\begin{equation}
 \boxed{\displaystyle
 \Delta C_{6,\mathrm{conn},3\TT}'=
 +\frac{7}{5120}\frac{g^4q^2b^6}{p^3a_*^4}L^{10}\,[1+o(1)].}
 \label{eq:onespecies}
\end{equation}
Here \(o(1)\to0\) as \(q/b\), \(bL\) and the relative scale-factor variation go to zero; the sign holds at small nonzero finite ratios. The factor is \((1/4)(2)(189/32)/(10\cdot6^3)\). The different signs in \eqref{eq:three} and \eqref{eq:onespecies} arise from polarization, not a switch of comparison convention. The matching of complete mixed sectors remains required for this same-species test; the contact isolation below uses distinguishable species.

\subsection{Real apertures and weak material calibration}\label{sec:packets}
The momentum-conserving kernels are continuous near both examples. For the second example this applies to every pairing because no transfer is zero. Smooth nonnegative momentum profiles concentrated about the six selected nonzero centers therefore integrate the signed primed kernel against a positive common momentum-conservation measure and preserve nonzero output. The unprimed result is the smeared version of the delta-normalized expressions, not a point evaluation of a plane-wave distribution. Each complex profile is the sum of a real aperture and \(\ii\) times a real aperture in physical position space. By multilinearity of the connected coefficient, at least one finite product of real records is nonzero. On the regular material-clock gauge slice these are configuration-field observables at one equal-clock time; the coordinates can be fixed by configuration-only gauge conditions, leaving their equal-slice configuration brackets unchanged. This is not a claim of general spacelike commutativity of nonlinear relational observables.

Smooth compact momentum profiles have rapidly decreasing spatial profiles. If actual compact spatial support is wanted, smooth spatial truncation converges to them in Schwartz seminorms and preserves a nonzero result, provided the same tempered regulated calculation is used for the truncations. Exact simultaneous compact support in momentum and position is neither used nor possible. A strict test band is imposed by momentum smearing; a strictly compact spatial detector has tails in momentum. One may make the packet hierarchy in a subcurvature effective theory, for example \(\Lambda_{\rm EFT}^{-1}\ll a_*L\ll a_*/b\), \(q\ll b\), \(a_*/q\ll R_{\mathrm{ap}}\ll H^{-1}\), choosing \(q\gg a_* H\) and a sufficiently weakly expanding patch. These are finite scale choices, not a claim the displayed effect is large. Products defining connected moments can be read on one equal-clock slice; finite narrow time gates approximate the smeared coefficients by continuity when taken as separate regulated correlation readouts.

The exact material-ruler pure TT action is
\(-\frac12\beta^2\int a^2\tr e^{-2gh}\).
Its quadratic term adds \(2\beta^2g^2 a^2\) to \(D_q\); its cubic is \((2/3)\beta^2g^3\int a^2\tr h^3\). The homogeneous clock contributes no pure tensor term through the determinant. There is no linear mixing of TT and minimally coupled reference perturbations: the one spatial gradient term pairs with the divergence of \(h\). For each equilateral \(M_b\) permutation \(\tr(M_bM_cM_d)=-9/8\), since the planar triple trace is \(-1/8\) and the normal triple trace is \(-1\). Thus the coefficient in \eqref{eq:three} multiplies by \(1-4\beta^2g^2/(7q^2)\), with the common modified retarded inverse throughout, if this additional pure TT cubic is truncated in the same way. The sign and the other finite pure-cubic witness persist for sufficiently weak nonzero ruler stress by continuity; one may instead compare the separated Einstein contribution holding the material one fixed. No material or graviton fluctuations are made new calibrated degrees of freedom to fit either arm.

\section{Distinguishing the bulk term from local mixed omissions}\label{sec:null}

This section uses three distinguishable scalar probes. Keep all propagating modes and entire apparatus, initial and endpoint sectors common. At the same tree order allow a specified set of additional omitted time-local two-tensor/two-scalar contour contacts of the form \(h_a h_a Q_a\), including a finite number of derivatives. This includes the missing mixed term in the spatial contact of Section~\ref{sec:cubic}. We do not claim a contact parametrization of the entire discrepancy for an arbitrary deterministic nonlinear metric.

\subsection{Local contact functions and bulk branches}
At this order, a connected difference from such a vertex has precisely one mixed insertion and two linear current insertions. The three species with two external fields each fix the possible attachments. Contracting the two tensor difference fields with the two other species gives advanced response legs, possibly differentiated at the vertex. The species whose fields occur inside \(Q_a\) instead gives an unpropagated local free scalar-pair factor. These properties follow from the \(aa\) contraction being zero. Derivatives can first be taken as one-sided interior time jets, with all boundary terms from integrations by parts retained and matched separately. We do not differentiate an unmatched hard switching distribution.

For each species use the routing of \eqref{eq:inplanesplit}, now with its own \(p_b>q/2\) and \(r_b=\sqrt{p_b^2-q^2/4}\). Keep \(q\), the three nonzero transfer vectors and their chosen planar polarizations fixed, and vary the three pair energies independently in open finite test bands. A local pair factor at \(v<T\) falls in a finite real span \(\mathcal V_b(v)\) of functions
\begin{equation}
 F_{b,\alpha}(p_b;v)
 =A_{b,\alpha}(p_b;v)\sin[2p_b(T-v)]
  +B_{b,\alpha}(p_b;v)\cos[2p_b(T-v)] ,
 \label{eq:contact}
\end{equation}
where \(A,B\) are algebraic in \(p_b\) at each fixed \(v\) on the admitted branch. Both factors from external free conformal scalar lines carry phases at the same vertex time; each local derivative multiplies a line by its energy or spatial momentum, apart from time-dependent background coefficients. Instantaneous rational spatial symbols in the derivative truncation are allowed if regular on the bands. They can depend on the scalar pair's individual spatial momenta and the other two \emph{fixed} tensor transfer vectors, not on the other two pair energies except through the already contracted advanced tensor legs. This split is an assumption on the contact family and licenses the independent controls. It would fail for an unmatched history-dependent detector or a genuinely dynamic exchanged reference field. Arbitrary spatial pseudodifferential symbols with no specified analytic dispersion are not implied.

The corresponding propagated bulk factor in \eqref{eq:three}, ignoring the common \(g\) and fixed polarization, is
\begin{equation}
 D_b(p_b;v)=(p_b^2-q^2/4)\int_v^T\dd\eta\,
          G_R(q;\eta,v)\frac{\sin[2p_b(T-\eta)]}{2p_b^2}.
 \label{eq:db}
\end{equation}
The notation \(D_b(p_b;v)\) here is a scalar form factor, not the wave operator \(D_q\).
\begin{proposition}[Local--bulk separation]
\label{prop:nonspan}
If \(G_R(q;T,v)\ne0\), the function \eqref{eq:db} is not in \(\mathcal V_b(v)\) on any nonempty regular open test band.
\end{proposition}
\begin{proof}
For the finite-derivative and rational spatial symbols above, fix the regular branch and continue along the positive real \(p_b\) end, across isolated poles by analytic continuation away from them if necessary. More generally restrict algebraic symbols to those admitting such continuation; after going round any intermediate singularities, each branch on the real end has a unique Puiseux asymptotic. The left side \eqref{eq:db} itself is analytic in \(p_b>q/2\). Twice integrating by parts with respect to \(\eta\), and using \(G_R(q;v,v)=0\), gives
\begin{equation}
 D_b(p_b;v)=\frac{G_R(q;T,v)}{4p_b}+O(p_b^{-2}).
 \label{eq:continuation}
\end{equation}
The upper endpoint contributes the first displayed term; the once differentiated sine boundary and the remaining integral are bounded after the second integration. If the alleged identity with \eqref{eq:contact} held on the band, it would hold throughout the end on those continued branches. On two sequences with \(\sin[2p_b(T-v)]=0\) and cosine \(+1\) and \(-1\), respectively, \(p_bB(p_b)\) would then approach the two distinct constants \(G_R(q;T,v)/4\) and \(-G_R(q;T,v)/4\). A single algebraic branch cannot do this. For flat \(a=1\) an elementary convolution gives the exact check
\[
 D_b(p_b;v)=\frac{(2p_b/q)\sin[q(T-v)]-\sin[2p_b(T-v)]}{8p_b^2}.
\]
The large-\(p_b\) step proves independence of finite-band kernels; it is not an extrapolation of a physical effective theory to measured large energy.
\end{proof}

\subsection{Start-time subtraction and finite tests}
Fix a lower slice \(v\) with the nonzero short-window coefficient above. For each \(b\), pick a basis of \(\mathcal V_b(v)\), append \(D_b(\cdot;v)\), and take as many evaluations as functions on its regular open band so that the evaluation matrix is nonsingular. Such evaluation points exist: otherwise the evaluation vectors fail to span, giving a linear identity. The real weights solving the associated linear system give
\begin{equation}
 \sum_m w_{b,m}F_{b,\alpha}(p_{b,m};v)=0\quad\text{for all }\alpha,
 \qquad
 \sum_m w_{b,m}D_b(p_{b,m};v)\ne0 .
 \label{eq:weights}
\end{equation}
The same weights are used across preparations. A pair energy refers to a choice of finite test aperture, not a change of free state or of the tensor spatial routing.

For a start slice \(s\) in an open neighborhood of \(v\) use the restrictions of the same free tensor Gaussian state and the same free conformal state. At this order each unmatched interior vertex has a single integral from \(s\) to \(T\). The advanced legs start at their vertex, not at \(s\), and the scalar contractions and common boundary prescription are independent of this variation. Thus the derivative in \(s\) sets the unmatched vertex time to \(s\), with a minus sign. Use the product weights \(\prod_b w_{b,m_b}\) across the finitely many triples of pair centers. At \(s=v\) every contact derivative contains one of the annihilated factors \eqref{eq:weights}, while the bulk derivative is
\begin{equation}
 \tfrac{63}{32}\,g^4 q^2 a(v)^2
   \prod_{b=1}^3\left[\sum_m w_{b,m}D_b(p_{b,m};v)\right]\ne0 .
 \label{eq:nullsignal}
\end{equation}
The signs follow from differentiating the lower limit in \eqref{eq:three}. The weak-ruler version uses an additional known multiplicative \(1-4\beta^2g^2/(7q^2)\ne0\) and the common modified Green function. This derivative statement is exact for the given allowed finite jet structures with arbitrary smooth coefficients, as long as the same coefficient profiles and state restrictions are used at all starts.

For every \emph{fixed} such smooth nuisance prescription, the two-run finite difference at \(v\) and \(v+\epsilon\) divided by \(\epsilon\) approaches the same nonzero result as \(\epsilon\to0\) for the ideal modes. After choosing a sufficiently small positive \(\epsilon\), choose the finitely many smooth apertures sufficiently narrow to preserve the nonzero difference. The narrow-aperture limit can first be calibrated by normalizing to the known positive common delta weight at each center tuple. A product of complex-mode profiles is a finite linear combination of real-mode data as in Section~\ref{sec:packets}; real and imaginary parts are real linear sums of actual measurements. This yields a finite nonzero combination, not a uniform finite tolerance over arbitrarily large or rapid contact coefficients, and its imperfect finite null on generic nuisance profiles is not called exact.

\begin{theorem}[An exact finite null for a finite list of templates]\label{thm:actualnull}
Fix a finite list of complete allowed mixed terms, including their time and background profiles, with arbitrary real amplitudes constant across the controls. Among the real connected six-point measurements with finite apertures and start slices in an open admissible interval there is a finite weighted combination that annihilates every term in the list exactly, but has nonzero pure bulk output. It need not use just two distinct start slices.
\end{theorem}
\begin{proof}
View the pure response \(P\) and each template response \(M_i\) as real functions on actual settings: six finite real spatial profiles, a real common final time as above, and the starting slice. Include the choice of running at either of two slices against a fixed common baseline if only differences are desired. If \(P\) were a real linear combination of the \(M_i\) on all settings, multilinearity would extend the identity to complex profiles. For each fixed nonzero-transfer center tuple, concentrate admissible smooth momentum profiles about it and divide by the common positive conservation-delta weight. The continuous primed kernels would satisfy the identical linear combination at those centers for all start slices. Differentiating that identity in the start slice at \(v\) gives zero on the mixed side under \eqref{eq:weights}, but gives \eqref{eq:nullsignal} on the pure side. This contradiction proves \(P\) independent of the finite span on the actual settings. Choosing an evaluation matrix for a basis of that span together with \(P\) now gives finite real weights on the actual settings (at most one more than the span dimension) with the stated exact outcome. If exactly compact spatial profiles rather than momentum-supported profiles are taken as admissible, the same implication uses their approximation in the common tempered calculation as in Section~\ref{sec:packets} before the momentum concentration. It does not require that a compact profile have a sharp spectral cutoff.
\end{proof}

The theorem does not posit fixed coefficients for arbitrary unknown time profiles. Prescribing the entire templates makes their global span finite; changing any coefficient as a function of the control would instead change the problem. Nor does the theorem isolate a bulk response from unmatched nonlocal mixed channels or a changed boundary functional.


\section{Relation to previous results and limitations}
A nonzero tensor-exchange scalar trispectrum in inflation is established \cite{SeerySlothVernizzi}; exchange wavefunctions and their contact ambiguities also have independent constructions \cite{JPS}. Equation~\eqref{eq:gauss} compares entire finite-time Gaussian influences with canonical covariance and response held fixed, not just one exchange limit. Ballesteros, Gamb\'in Egea and P\'erez Rodr\'iguez \cite{Ballesteros} analyze differences between quantum and classical evolution from a finite inflationary initial time, including scalar bispectrum and tensor power examples; a flattened pole is not needed. Hauser and Laine \cite{HauserLaine} keep a quantum scalar source in a classical retarded calculation and compare second-order and one-loop tensor power. In stochastic semiclassical gravity, Hu, Roura and Verdaguer \cite{HuRouraVerdaguer} relate an Einstein--Langevin equation to symmetrized quantum metric correlations at leading large \(N\), rather than to the complete doubled scalar response used here. The separated nonlinear signal in this paper concerns the tensor self-cubic and an external connected scalar six-point coefficient. Its nonzero realization and finite mixed-response test, not the general existence of quantum response vertices, are the additional statements.

The result neither establishes an autonomous completely positive classical comparator nor rules out engineered non-Gaussian tensor noise, uncontrolled time-nonlocal interactions, other propagating modes, or different detector calibrations. It does not identify an unprocessed primordial sky observable, an entanglement witness, or the first quantum correction to every correlator. No higher-loop or ultraviolet completion is used.

\begin{thebibliography}{9}
\bibitem{FeynmanVernon}
R.~P. Feynman and F.~L. Vernon, Jr.,
The theory of a general quantum system interacting with a linear dissipative system,
\emph{Ann. Phys.} \textbf{24}, 118--173 (1963),
\href{https://doi.org/10.1016/0003-4916(63)90068-X}{doi:10.1016/0003-4916(63)90068-X}.
\bibitem{Gao}
X. Gao et al.,
Full bispectra from primordial scalar and tensor perturbations in the most general single-field inflation model,
\emph{PTEP} \textbf{2013}, 053E03 (2013),
\href{https://doi.org/10.1093/ptep/ptt031}{doi:10.1093/ptep/ptt031}.
\bibitem{Polkovnikov}
A. Polkovnikov, Phase space representation of quantum dynamics,
\emph{Ann. Phys.} \textbf{325}, 1790--1852 (2010),
\href{https://arxiv.org/abs/0905.3384}{arXiv:0905.3384}.
\bibitem{SeerySlothVernizzi}
D. Seery, M.~S. Sloth and F. Vernizzi, Inflationary trispectrum from graviton exchange,
\emph{JCAP} \textbf{03} (2009) 018,
\href{https://arxiv.org/abs/0811.3934}{arXiv:0811.3934}.
\bibitem{JPS}
S. Jazayeri, E. Pajer and D. Stefanyszyn, From locality and unitarity to cosmological correlators,
\emph{JHEP} \textbf{10} (2021) 065,
\href{https://arxiv.org/abs/2103.08649}{arXiv:2103.08649}.
\bibitem{Ballesteros}
G. Ballesteros, J. Gamb\'in Egea and A. P\'erez Rodr\'iguez,
Classical and quantum evolution of inflationary fluctuations,
\href{https://doi.org/10.1103/w7y3-gxbr}{doi:10.1103/w7y3-gxbr},
\href{https://arxiv.org/abs/2604.18416}{arXiv:2604.18416}.
\bibitem{HauserLaine}
A. Hauser and M. Laine,
Matching second-order classical and 1-loop quantum tensor power spectra in de Sitter spacetime,
\emph{JHEP} \textbf{11} (2026) 167,
\href{https://doi.org/10.1007/JHEP11(2026)167}{doi:10.1007/JHEP11(2026)167}.
\bibitem{HuRouraVerdaguer}
B.~L. Hu, A. Roura and E. Verdaguer,
Induced quantum metric fluctuations and the validity of semiclassical gravity,
\emph{Phys. Rev. D} \textbf{70}, 044002 (2004),
\href{https://arxiv.org/abs/gr-qc/0402029}{arXiv:gr-qc/0402029}.
\end{thebibliography}
\end{document}
